\section{Introduction}
\label{sec:introduction}

LLM-based agents increasingly interact with external tool ecosystems to accomplish real-world tasks \cite{yao2023react, shen2024hugginggpt}.
As these ecosystems grow from tens to hundreds of tools, \emph{tool routing}---determining which tools to invoke for a given query---becomes a critical bottleneck.
Direct tool selection, where the LLM chooses from the full catalog, degrades in accuracy, introduces hallucinated tool invocations, and scales poorly with catalog size \cite{li2025toolscope, qin2023toolllm}.

We argue that the fundamental issue is the \emph{abstraction}: framing tool routing as tool selection forces the model to simultaneously identify relevant tools, determine their ordering, and ensure compositional validity---all in a single unconstrained decision.

\begin{figure*}[t]
\centering
\begin{tikzpicture}[
    node distance=0.5cm,
    box/.style={rectangle, draw, rounded corners=2pt, minimum width=2.8cm, minimum height=0.6cm, font=\small, align=center},
    arrow/.style={-{Stealth[length=4pt]}, thick},
    annot/.style={font=\footnotesize\itshape, text=gray!70!black}
]

% Traditional approach (left)
\node[font=\small\bfseries] (trad-title) {Traditional Tool Routing};

\node[box, below=0.4cm of trad-title, fill=blue!8] (query1) {Query};
\node[box, below=0.5cm of query1, fill=orange!10, minimum width=3.2cm, minimum height=1.0cm] (llm1) {LLM\\{\scriptsize choose from 135 tools}};
\node[box, below=0.5cm of llm1, fill=red!8] (tools1) {Selected tools};

\draw[arrow] (query1) -- (llm1);
\draw[arrow] (llm1) -- (tools1);

\node[annot, below=0.15cm of tools1] (prob1) {hallucinations, missed steps, F1 $\approx$ 0.50};

% Divider
\draw[dashed, gray!50, line width=0.5pt] (4.2, 0.5) -- (4.2, -4.5);

% Our approach (right)
\node[font=\small\bfseries, right=5.5cm of trad-title] (our-title) {Typed Composition Search};

\node[box, below=0.4cm of our-title, fill=blue!8] (query2) {Query};
\node[box, below=0.5cm of query2, fill=green!10, minimum width=3.5cm] (types) {LLM: predict types\\{\scriptsize (Namespace, PodLogs)}};
\node[box, below=0.5cm of types, fill=yellow!8, minimum width=3.5cm] (graph) {Graph BFS\\{\scriptsize Namespace $\to$ Pod $\to$ PodLogs}};
\node[box, below=0.5cm of graph, fill=green!15] (tools2) {3 tools returned};

\draw[arrow] (query2) -- (types);
\draw[arrow] (types) -- (graph);
\draw[arrow] (graph) -- (tools2);

\node[annot, below=0.15cm of tools2] (prob2) {0 hallucinations, auto multi-hop, F1 $\approx$ 0.85};

\end{tikzpicture}
\caption{Approach comparison. \textbf{Left:} Traditional tool routing selects from the full catalog in a single unconstrained step. \textbf{Right:} Typed Composition Search predicts entity types, then recovers the tool sequence through graph search. The graph guarantees compositional validity and eliminates hallucinated tools.}
\label{fig:approach}
\end{figure*}

We propose \textbf{Typed Composition Search}, a reformulation that decomposes tool routing into \emph{entity type prediction} and \emph{graph-based path recovery} (Figure~\ref{fig:approach}).
Tools are modeled as typed transformations between entity types (e.g., \texttt{get\_pod\_logs}: \texttt{Pod} $\rightarrow$ \texttt{PodLogs}), forming a directed graph where entity types are nodes and tools are edges.
Given a query, the LLM predicts the source and target entity types; breadth-first search on the graph recovers the tool sequence.

This decomposition separates semantic reasoning (what entities does the user mean?) from compositional reasoning (what tools connect those entities?), assigning each to the mechanism best suited for it.

We evaluate this approach across four models (8B--20B parameters plus one proprietary), five tool domains (54--170 tools), and 140 queries in six difficulty categories.
Our contributions are:

\begin{enumerate}
    \item \textbf{Formulation.} We reformulate tool routing as typed entity reasoning, shifting the LLM's task from tool selection to entity classification.

    \item \textbf{Execution model.} We propose a typed graph execution model combining entity prediction with deterministic graph search. No training, fine-tuning, or embedding computation is required.

    \item \textbf{Analytical decomposition.} We show that end-to-end recall decomposes into type prediction accuracy and graph reachability (Eq.~\ref{eq:recall-decomposition}), enabling targeted failure analysis.

    \item \textbf{Benchmark.} We provide a comprehensive evaluation across 4 models, 5 domains, and 6 routing strategies, covering 140 queries in 6 difficulty categories.

    \item \textbf{Empirical evidence.} Graph routing consistently improves tool selection (19/19 model--domain combinations, avg.\ $\Delta$F1 = +0.33) while eliminating structurally invalid tool calls.
\end{enumerate}
